\section{The complete fixed-relative-accuracy staircase}\label{sec:fixed}
All approximation-order indices, including $r$ and $j$, are nonnegative integers. For $r\geq0$ define
\begin{equation}\label{eq:G-definition}
 G_r(\rho)=\inf\left\{\max_{1\leq y\leq\rho}|P(y)-y|:
 P\in\R[y],\ \deg P\leq2r,\ P\geq0\text{ on }\R\right\}.
\end{equation}
The constraint on the entire real line is essential: the approximation error is measured only on $[1,\rho]$.

\begin{theorem}[Optimal fixed-accuracy queries]\label{thm:staircase}
Fix $\rho>1$, $0<c\leq1$, and $K>0$, and put
$\ell=\min\{j\geq0:G_j(\rho)\leq K\}$.
Then $\ell$ is finite and, as $\delta\downarrow0$,
\begin{equation}\label{eq:staircase}
 Q(\delta,K\delta;\rho,c)=
 \begin{cases}
 \Theta_{\rho,c,K}(\log(1/\delta)),&\ell=0,\quad c<1,\\
 \Theta_{\rho,K}(1),&\ell=0,\quad c=1,\\
 \Theta_{\ell,\rho,c,K}(\delta^{-1+1/(2\ell)}),&\ell\geq1.
 \end{cases}
\end{equation}
Every threshold equality $K=G_\ell$ belongs to the cheaper tier displayed in this formula.
\end{theorem}
The distinction at $c=1$ has a simple cause: $[c,1]$ is then a single known spectral point. A high interval of positive length supplies an additional logarithmic obstruction, whereas one known point can be handled by interpolation.

\subsection{The globally nonnegative approximation problem}
Write
\[
 m_\rho=\frac{\rho+1}{2},\qquad h_\rho=\frac{\rho-1}{2},\qquad
 z_0=\frac{m_\rho}{h_\rho},
\]
and let $T_n$ denote the Chebyshev polynomial, $T_n(\cos\theta)=\cos(n\theta)$.
\begin{theorem}[Exact thresholds]\label{thm:thresholds}
We have $G_0=h_\rho$. For $r\geq1$,
\begin{align}
 G_r&=h_\rho\max_{z\geq z_0}\frac{z-z_0}{T_{2r}(z)},\label{eq:G-formula}\\
 P_r^*(y)&=y+G_r T_{2r}\bigl((y-m_\rho)/h_\rho\bigr)\label{eq:extremizer}
\end{align}
are respectively the optimal error and an admissible minimizer. The positive numbers $G_r$ strictly decrease to zero. If $z_*$ is the maximizer in \eqref{eq:G-formula}, it is unique and
\begin{equation}\label{eq:G-stationary}
 z_*-z_0=T_{2r}(z_*)/T_{2r}'(z_*),\qquad
 G_r=h_\rho/T_{2r}'(z_*).
\end{equation}
\end{theorem}
\begin{proof}
The constant case is ordinary interval minimax approximation. For $n=2r\geq2$, the exterior Chebyshev inequality says that a degree-at-most-$n$ polynomial $A$ bounded by one on $[-1,1]$ satisfies $|A(z)|\leq T_n(|z|)$ for $|z|>1$. One elementary proof compares $A$ with $T_n$ at its $n+1$ successive extrema: if the exterior inequality failed, subtracting a suitable multiple of $T_n$ would give $n$ interior sign changes and an exterior zero, contradicting the degree. Apply this inequality to $P(y)-y$. At $y=-h_\rho(z-z_0)\leq0$, nonnegativity gives $P(y)-y\geq h_\rho(z-z_0)$, proving the lower bound in \eqref{eq:G-formula}.

To prove admissibility of \eqref{eq:extremizer}, set $e=G_r/h_\rho$ and $x=(y-m_\rho)/h_\rho$. We must show $x+z_0+eT_n(x)\geq0$ for every real $x$. For $x\leq-z_0$ this is precisely the maximizing inequality defining $e$. For $-z_0\leq x\leq-1$ and for $x\geq1$, both terms are nonnegative. On $[-1,1]$, only negative lobes of $T_n$ require attention. They satisfy $x\geq-\cos(\pi/(2n))$, so
\[
 x+z_0>1-\cos(\pi/(2n))\geq n^{-2}.
\]
The last bound follows from $\sin u\geq2\sqrt2 u/\pi$ on $[0,\pi/4]$. Since $T_n$ is strictly convex on $[1,\infty)$ and $T_n'(1)=n^2$,
$T_n(z)\geq1+n^2(z-1)$, whence $e<n^{-2}$. Since $T_n(x)\geq-1$, positivity follows. The approximation error is exactly $G_r$ by Chebyshev alternation.

The maximum in \eqref{eq:G-formula} is positive, attained away from $z_0$ and infinity. The function $z-T_n(z)/T_n'(z)$ increases strictly on $(1,\infty)$, with derivative $T_nT_n''/(T_n')^2>0$, and tends to infinity. This proves uniqueness and \eqref{eq:G-stationary}. For successive positive indices, evaluate the ratio for $T_{2r+2}$ at its maximizer and use $T_{2r+2}(z)>T_{2r}(z)$ for $z>1$. Also $G_1<G_0$: the explicit formula below proves this directly. Finally, polynomial approximations to $\sqrt y$ on $[1,\rho]$, squared, are globally nonnegative and converge uniformly to $y$. Thus $G_r\to0$.
\end{proof}
The admissibility has a geometric antecedent: by \eqref{eq:G-stationary}, in the coordinate $z=(m_\rho-y)/h_\rho$, the extremizer is $G_r$ times the difference between $T_{2r}(z)$ and its tangent at $z_*$. Its nonnegativity also follows from the supporting-tangent inequality of Fawzi, Saunderson, and Parrilo \cite[Appendix A, Lemma 1]{fawzi2017}; the proof above is self-contained.

For the first threshold the formula simplifies to
\begin{equation}\label{eq:G1}
 G_1=\frac{\rho+1-\sqrt{(\rho^2+6\rho+1)/2}}4,
 \qquad P_1^*(y)=\frac{(y+s)^2}{2(m_\rho+s)},\quad
 s=\sqrt{m_\rho^2-h_\rho^2/2}.
\end{equation}
These expressions follow either by solving \eqref{eq:G-stationary} for $T_2(z)=2z^2-1$ or by expanding \eqref{eq:extremizer}.

\begin{proposition}[Threshold asymptotics]\label{prop:threshold-asymptotics}
Let $R_0=(\sqrt\rho+1)/(\sqrt\rho-1)$ and $a=\log R_0$. As $r\to\infty$,
\begin{equation}\label{eq:G-asymptotic}
 G_r=\frac{\sqrt\rho}{\mathrm e\,r}R_0^{-2r}
 (1+O_\rho(r^{-1})).
\end{equation}
Consequently, as $K\downarrow0$, with $L=\log(1/K)$,
\[
 \ell(K)=\frac{L-\log L+\log(2\sqrt\rho\,a/\mathrm e)}{2a}+O_\rho(1),
 \qquad
 1-\frac1{2\ell(K)}
 =1-\frac{a}{L-\log L+O_\rho(1)}.
\]
The $O(1)$ term in the first expression includes integer rounding.
\end{proposition}
\begin{proof}
Put $n=2r$, $z=\cosh t$, and note $z_0=\cosh a$. The stationary equation is
\[
 n(\cosh t-\cosh a)\tanh(nt)=\sinh t.
\]
Its solution has $t=a+O_\rho(n^{-1})$: indeed the increasing left-to-right ratio $n(\cosh t-\cosh a)/\sinh t$ crosses one within such a neighborhood, and $\tanh(nt)=1+O_\rho(e^{-2na})$. Expanding with $t=a+s/n$ now gives $s=1+O_\rho(n^{-1})$. Substitution into \eqref{eq:G-formula}, using $h_\rho\sinh a=\sqrt\rho$, gives \eqref{eq:G-asymptotic}. Taking logarithms and inverting $L=2ar+\log r-\log(\sqrt\rho/\mathrm e)+O_\rho(1/r)$ gives the last claims, with bounded rounding uncertainty.
\end{proof}

\subsection{An all-circuit lower bound}
\begin{theorem}[Nonnegative Taylor limit]\label{thm:hierarchy}
Fix $r\geq0$ and $K<G_r(\rho)$. Any coherent converter accurate to $K\delta$ on $[\delta,\rho\delta]$ uses
\[
 T=\Omega_{r,\rho,K}(\delta^{-(2r+1)/(2r+2)})
\]
queries. No upper spectral interval is needed for this bound.
\end{theorem}
\begin{proof}
Use Lemma~\ref{lem:scalar} and put $g_\delta(x)=1-\operatorname{Re}q_\delta(\arcsin x)$. Correctness implies
\[
 |g_\delta(\delta y)/\delta-y|\leq K\qquad(1\leq y\leq\rho).
\]
Taylor expansion through order $2r+1$, followed by rescaling, gives a polynomial $A_\delta$ of degree at most $2r+1$ such that, on each fixed compact set of real $y$,
\[
 g_\delta(\delta y)/\delta=A_\delta(y)
 +O_{r,y}((T+1)^{2r+2}\delta^{2r+1}).
\]
If the claimed bound failed, there would be a sequence $\delta\downarrow0$ with $T\delta^{(2r+1)/(2r+2)}\to0$, so the remainder would vanish uniformly on each fixed compact set. Values at $2r+2$ fixed distinct points of $(1,\rho)$ then bound all coefficients of $A_\delta$ by invertibility of the Vandermonde matrix. Pass to a coefficientwise convergent subsequence, with limit $A$. For each fixed real $y$, global nonnegativity of $g_\delta$ gives $A(y)\geq0$. A nonzero nonnegative polynomial has even degree; thus $\deg A\leq2r$. It approximates $y$ within $K$ on $[1,\rho]$, contradicting the definition of $G_r$.
\end{proof}
If $\eta(\delta)/\delta\to0$, this immediately yields $T=\Omega_{\rho,\varepsilon}(\delta^{-1+\varepsilon})$ for every fixed $\varepsilon>0$ (with the onset depending on the accuracy function). This is stronger than the pairwise square-root obstruction. It uses contractivity outside the correctness interval, through the globally nonnegative Taylor limit.

\subsection{Bounded constructions, including exact threshold errors}
We first record a positive approximant for the fixed upper interval. Define $d_j$ by
\[
 \frac{1-\sqrt{1-z}}{1-z}=\sum_{j\geq1}d_jz^j,\qquad
 S_N(u)=u\sum_{j=1}^N d_j(1-u)^j.
\]
The coefficients $d_j$ are partial sums of the positive Taylor coefficients of $1-\sqrt{1-z}$, so $0<d_j\leq1$. For $0\leq u\leq1$,
\begin{equation}\label{eq:SN}
 0\leq S_N(u)\leq1-\sqrt u,\quad
 0\leq1-\sqrt u-S_N(u)\leq(1-u)^{N+1},\quad S_N(u)\leq Nu.
\end{equation}
The tail estimate follows by comparison with the geometric series; the endpoints follow by continuity. For $c<1$, $N=O_c(\log(1/\delta))$ makes the high-band tail any prescribed constant multiple of $\delta$. For $c=1$, $N=1$ suffices because $S_N(1)=0$.

\begin{lemma}[A pinned gate]\label{lem:pinned}
Fix $r\geq1$ and a nonempty finite set $A\subseteq[1,\rho]$. Put $\alpha=1-1/(2r)$ and choose an odd integer $m\asymp M\delta^{-\alpha}$, where $M>0$ is fixed. There is an even polynomial $W$ of degree $O_{r,|A|}(m)$ with $0\leq W\leq1$ on $[-1,1]$, $W(\delta y)=1$ for $y\in A$, and
\begin{align}
 1-W(\delta y)&\leq C(m\delta)^{2r+2}\dist(y,A)^{2r+2},
 &&1\leq y\leq\rho,\label{eq:pinned-low}\\
 W(x)&\leq C(m|x|)^{-2r-4},&& C_0\delta\leq|x|\leq1.\label{eq:pinned-tail}
\end{align}
Constants may depend on $r,\rho,A$, but not on $\delta,m$; $C_0$ is chosen sufficiently large depending on $\rho$.
\end{lemma}
\begin{proof}
For $\theta_j=\arcsin(\delta y_j)$, $y_j\in A$, define
\[
 D_{m,a}(\theta)=\frac{\sin(m(\theta-a))}{m\sin(\theta-a)},\quad
 W(\sin\theta)=1-
 \prod_{y_j\in A}\prod_{\sigma=\pm1}
 \left(1-D_{m,\sigma\theta_j}(\theta)^{2r+4}\right)^{r+1}.
\]
Removable singularities are filled by continuity. The elementary inequality $|\sin(mu)|\leq m|\sin u|$ gives $|D_{m,a}|\leq1$. For odd $m$, $D_{m,a}$ is a real trigonometric polynomial with even frequencies and degree $m-1$. Pairing $a$ with $-a$ makes the product even in $\theta$ and $\pi$-periodic. It is therefore a polynomial in $\cos2\theta=1-2x^2$, of degree $O_{r,|A|}(m)$ in $x$. The range and pinning properties are immediate.

Near a pin, $1-D_{m,a}(\theta)^{2r+4}\leq C_rm^2(\theta-a)^2$. Select the nearest positive pin in the product defining $1-W$ and bound all other factors by one. Since $|\arcsin(\delta y)-\arcsin(\delta y_j)|\leq2\delta|y-y_j|$ for small $\delta$, this proves \eqref{eq:pinned-low}. Finally, if $C_0\delta\leq|x|\leq1$, then $|\sin(\arcsin x\pm\theta_j)|\geq|x|/2$. Hence $|D_{m,\pm\theta_j}|\leq2/(m|x|)$. Applying $1-\prod_i(1-a_i)^{r+1}\leq(r+1)\sum_i a_i$ proves \eqref{eq:pinned-tail}.
\end{proof}

\begin{proposition}[Upper bound at every positive threshold]\label{prop:threshold-upper}
For each fixed $r\geq1$ there is a polynomial $q_\delta$ bounded by one on $[-1,1]$ with error at most $G_r\delta$ on \eqref{eq:bands} and degree $O_{r,\rho,c}(\delta^{-1+1/(2r)})$.
\end{proposition}
\begin{proof}
Take $P=P_r^*$ and $e(y)=P(y)-y$. The positive-contact set
$A=\{y\in[1,\rho]:e(y)=G_r\}$ is finite and nonempty, by \eqref{eq:extremizer}. Use Lemma~\ref{lem:pinned} for this set and define
\begin{equation}\label{eq:blend}
 q_\delta(x)=W(x)[1-\delta P(x/\delta)]+(1-W(x))S_N(x^2).
\end{equation}
We verify contractivity before its approximation property. Write $q_\delta=B-\delta P(x/\delta)W$ with $B=W+(1-W)S_N(x^2)\in[0,1]$. Thus $q_\delta\leq1$. Since $P(y)\leq C(1+|y|^{2r})$ and $P\geq0$, splitting at $|x|=C_0\delta$ and $|x|=1/m$ gives
\begin{equation}\label{eq:global-loss}
 0\leq\delta P(x/\delta)W(x)
 \leq C\delta+C\delta^{1-2r}m^{-2r}
 \leq C\delta+C'M^{-2r}.
\end{equation}
Indeed $C_0\delta<1/m$ for small $\delta$; on the middle region use $W\leq1$, and on the outer region use \eqref{eq:pinned-tail}. Choose $M$ large enough and then $\delta$ small enough to make \eqref{eq:global-loss} at most one. It follows that $q_\delta\geq-1$.

For the lower interval, the polynomial $G_r-e(y)$ is nonnegative on $[1,\rho]$ and has zeros there exactly at $A$, with multiplicities at most $2r$. Compactness and factorization therefore give a constant $b>0$ such that
\[
 G_r-e(y)\geq b\dist(y,A)^{2r}\quad(1\leq y\leq\rho).
\]
Since $(m\delta)^{2r+2}=O(\delta^{1+1/r})$, \eqref{eq:pinned-low} and $W=1-o(1)$ on this interval imply, for small $\delta$,
\begin{equation}\label{eq:pinned-slack}
 1-W(\delta y)\leq\tfrac12 W(\delta y)\delta[G_r-e(y)].
\end{equation}
This estimate includes the contacts themselves, where both sides vanish. Put $d(x)=1-x-S_N(x^2)\in[0,1]$ for $0\leq x\leq1$. On $x=\delta y$,
\[
 (1-x)-q_\delta(x)=W\delta e(y)+(1-W)d(x).
\]
Its lower bound is $-G_r\delta$ because $e\geq-G_r$ and $d\geq0$. Its upper bound follows from \eqref{eq:pinned-slack}:
\[
 G_r\delta-[(1-x)-q_\delta(x)]
 =W\delta(G_r-e)+(1-W)(G_r\delta-d)\geq0.
\]
Thus the error is at most exactly $G_r\delta$, without an asymptotic error overshoot.

On $[c,1]$, the tail in \eqref{eq:SN} can be made at most $G_r\delta/3$. Also $W=O(m^{-2r-4})=o(\delta)$, and
\[
 \delta P(x/\delta)W=O(\delta^{1-2r}m^{-2r-4})
 =O(\delta^{4-2/r})=o(\delta).
\]
These two terms can each be made at most $G_r\delta/3$ by decreasing $\delta$. Finally, \eqref{eq:blend} has degree $O_{r,\rho}(m+N)$, proving the claim.
\end{proof}
The pinned construction handles strict bands as well, by using the smallest $r$ with $G_r\leq K$. Pinning is necessary in this proof to control the sign of the error at the exact minimax contacts; a gate with merely $o(\delta)$ leakage would only establish errors $(G_r+o(1))\delta$.

\subsection{The coarse tier and completion of the staircase proof}
\begin{proposition}[Coarse conversion]\label{prop:coarse}
For $K\geq G_0=(\rho-1)/2$, the first two lines of \eqref{eq:staircase} hold.
\end{proposition}
\begin{proof}
It suffices to construct error $G_0\delta$. Set $P=m_\rho$, so $e(y)=m_\rho-y$ and $G_0-e(y)=y-1$. First suppose $c<1$ and use
\[
 W(x)=[1-(x^2-\delta^2)^2]^L,\qquad
 L=\lceil C_c\log(1/\delta)\rceil.
\]
For small $\delta$, the base is in $[0,1]$ on $[-1,1]$. On the lower band,
\[
 1-W(\delta y)\leq L\delta^4(y^2-1)^2
 \leq C_\rho L\delta^4(y-1)^2.
\]
This is at most $\tfrac12W\delta(y-1)$ for sufficiently small $\delta$. On the high band, $W\leq(1-c^4/4)^L$ for $\delta$ sufficiently small. Choose the constant in $L$ so that this leakage is at most $G_0\delta/3$, and choose $N=O_c(\log(1/\delta))$ so the tail in \eqref{eq:SN} is at most $G_0\delta/3$. The blend \eqref{eq:blend} is a convex combination of $1-m_\rho\delta$ and $S_N(x^2)$, hence lies in $[0,1]$. The same signed-error argument as in Proposition~\ref{prop:threshold-upper} proves the exact low-band error. On the high band its error is at most the tail plus $W$, which is below $G_0\delta$. Its degree is $O(\log(1/\delta))$.

If $c=1$, instead take
\[
 W(x)=1-\left(\frac{x^2-\delta^2}{1-\delta^2}\right)^2,
 \qquad q_\delta(x)=(1-m_\rho\delta)W(x).
\]
For small $\delta$, $0\leq W\leq1$ on $[-1,1]$, $W(\delta)=1$, and $W(1)=0$. The low leakage is $O_\rho(\delta^4(y-1)^2)$ and the same signed-error argument applies with $S_N=0$. This degree-four polynomial is exact at the high point $1$. A zero-query circuit has the same scalar output at $x=\delta$ and $x=1$, whose targets differ by $1-\delta$, so cannot satisfy the error contract for small $\delta$. This gives $\Theta(1)$.

For the logarithmic lower bound when $c<1$, choose a closed interval $I$ of positive length $b$ strictly inside $(\arcsin c,\pi/2)$. Let $q(t)$ be the scalar output entry and $R(t)=q(t)-(1-\sin t)$. It has trigonometric degree $n\leq\max(T,1)$, is bounded by $K\delta$ on $I$, and satisfies $|R(-\pi/2)|\geq1$ by contractivity. Here is a direct form of the needed interpolation inequality. If $R$ has degree at most $n$, then $z^nR(t)$, $z=e^{it}$, is an ordinary polynomial of degree at most $2n$. Interpolate at $2n+1$ equally spaced angles in $I$. For node indices $i,j$, their distance on the unit circle is at least $b|i-j|/(\pi n)$, whereas its distance to any evaluation point on the unit circle is at most two. The sum of the absolute Lagrange weights is consequently at most
\[
 \left(\frac{2\pi n}{b}\right)^{2n}
 \sum_{j=0}^{2n}\frac1{j!(2n-j)!}
 =\frac{(4\pi n/b)^{2n}}{(2n)!}
 \leq(2\pi\mathrm e/b)^{2n}.
\]
Thus $1\leq K\delta(2\pi\mathrm e/b)^{2n}$, giving $T=\Omega_{c,K}(\log(1/\delta))$.
\end{proof}
\begin{proof}[Proof of Theorem~\ref{thm:staircase}]
Theorem~\ref{thm:thresholds} guarantees existence of $\ell$. For $\ell\geq1$, the lower bound is Theorem~\ref{thm:hierarchy} with $r=\ell-1$, since $K<G_{\ell-1}$. Proposition~\ref{prop:threshold-upper}, followed by Lemma~\ref{lem:implementation}, gives the matching upper bound because $G_\ell\leq K$. Proposition~\ref{prop:coarse} treats $\ell=0$.
\end{proof}

\subsection{Definite parity and a polynomial query separation}
Even polynomials satisfy additional lower bounds governed by the following approximation errors. Define
\[
 E_r(\rho)=\inf_{A\in\R[v],\,\deg A\leq r}
 \max_{1\leq y\leq\rho}|y-A(y^2)|.
\]
These unconstrained errors first give elementary lower-bound cutoffs. Retaining global nonnegativity below will give an optimal implicit even-parity staircase, including equality cases. Each $E_r$ is positive: otherwise finite-dimensional compactness would give an exact polynomial representation of $\sqrt v$ on $[1,\rho^2]$, which is impossible. In particular, the affine cutoff is
\begin{equation}\label{eq:E1}
 E_1=\frac{(\rho-1)^2}{8(\rho+1)},\qquad
 A_1(v)=a_\rho v+b_\rho,\quad
 a_\rho=\frac1{\rho+1},\quad
 b_\rho=\frac\rho{\rho+1}+E_1.
\end{equation}
To verify minimax optimality, the error $\sqrt v-A_1(v)$ equals $-E_1$ at $v=1,\rho^2$, and $E_1$ at $v=(\rho+1)^2/4$. An affine perturbation reducing all three absolute errors would need two sign changes, which is impossible.

\begin{proposition}[Parity obstructions]\label{prop:parity-lower}
Let $p$ be a real polynomial bounded by one on $[-1,1]$ that approximates $1-x$ within $K\delta$ on $[\delta,\rho\delta]$. If $p$ is even and $K<E_r(\rho)$, then
\[
 \deg p=\Omega_{r,\rho,K}(\delta^{-1+1/(2r+2)}).
\]
More explicitly, if $\deg p=2m$ and $K<E_1$, then
\begin{equation}\label{eq:even-explicit}
 m\geq\left(\frac{3(E_1-K)}{2\rho^4}\right)^{1/4}\delta^{-3/4}.
\end{equation}
If $p$ is odd, then
\[
 \deg p\geq\frac{[1-(1+K)\delta]\sqrt{1-\delta^2}}{\delta}.
\]
\end{proposition}
\begin{proof}
Taylor-expand an even degree-$d$ polynomial at zero through degree $2r$. On a fixed interior subinterval, Bernstein inequalities give $|p^{(j)}|\leq C_j(d+1)^j$ for each fixed $j$. The vanishing odd Taylor terms imply that $(1-p(\delta y))/\delta$ differs from a degree-$r$ polynomial in $y^2$ by at most $C_{r,\rho}(d+1)^{2r+2}\delta^{2r+1}$. Therefore
$E_r\leq K+C_{r,\rho}(d+1)^{2r+2}\delta^{2r+1}$, giving the first claim.
For the explicit constant, write $p(x)=A(x^2)$ with $\deg A=m$ and $\norm{A}_{[0,1]}\leq1$. The second Markov inequality on $[0,1]$ gives $\norm{A''}\leq4m^2(m^2-1)/3$. The first-order Taylor remainder at $u=\delta^2v$, divided by $\delta$, is at most $2m^4\rho^4\delta^3/3$. Comparing to the best affine error proves \eqref{eq:even-explicit}. Finally, odd parity gives $p(0)=0$, so the mean-value theorem and Bernstein's first-derivative inequality give
$(1-(1+K)\delta)/\delta\leq\sup_{[0,\delta]}|p'|\leq d/\sqrt{1-\delta^2}$.
\end{proof}

\begin{proposition}[Even square-root construction]\label{prop:even-upper}
For every fixed $\gamma>0$, there is an even polynomial bounded by one on $[-1,1]$ with error at most $(E_1+\gamma)\delta$ on \eqref{eq:bands} and degree $O_{\rho,c,\gamma}(\delta^{-1/2})$.
\end{proposition}
\begin{proof}
For an integer $m\geq2$ define the normalized endpoint Fej\'er kernel
\[
 F_m(u)=\frac{1-T_m(1-2u)}{2m^2u},\quad F_m(0)=1.
\]
This is a polynomial of degree $m-1$. If $u=\sin^2(\theta/2)$, then
$F_m(u)=\sin^2(m\theta/2)/(m^2\sin^2(\theta/2))$, so
$0\leq F_m\leq1$ on $[0,1]$ and $F_m(u)\leq1/(m^2u)$ for $0<u\leq1$.
Taylor expansion at zero gives
\[
 F_m(u)=1-\frac{m^2-1}{3}u+O(m^4u^2).
\]
The remainder is uniform when $m^2u=o(1)$; it also follows from the second Markov bound and $\norm{F_m}_{[0,1]}\leq1$.
Choose an integer $m$ nearest the positive solution of
$2(m^2-1)\delta/3=a_\rho$ and set $W_m=F_m^2$. Integer rounding gives, uniformly on $1\leq v\leq\rho^2$,
\[
 1-W_m(\delta^2v)=a_\rho\delta v+O_\rho(\delta^{3/2}).
\]
Using \eqref{eq:SN} with a high-band tail $O(\delta^2)$, define
\[
 P_\delta(u)=(1-b_\rho\delta)W_m(u)+(1-W_m(u))S_N(u),
 \qquad p_\delta(x)=P_\delta(x^2).
\]
This is a convex combination of numbers in $[0,1]$. On the low band, $S_N(\delta^2v)=O(N\delta^2)$ and hence
\[
 1-P_\delta(\delta^2v)=\delta(b_\rho+a_\rho v)
 +O_{\rho,c}(\delta^{3/2}+\delta^2\log(1/\delta)).
\]
On the high band, $W_m(u)\leq1/(m^4u^2)=O_{\rho,c}(\delta^2)$, so $P_\delta(u)=1-\sqrt u+O_{\rho,c}(\delta^2)$. The degree is at most $2[2(m-1)+N+1]=O_{\rho,c}(\delta^{-1/2})$, and \eqref{eq:E1} finishes the proof.
\end{proof}
\begin{definition}[Nonnegative even thresholds]\label{def:even-thresholds}
For $r\geq0$ let
\[
 F_r(\rho)=\inf_{\substack{P\in\R[y]\text{ even},\ \deg P\leq2r\\
                          P\geq0\text{ on }\R}}
                   \max_{1\leq y\leq\rho}|P(y)-y|.
\]
\end{definition}
These constrained thresholds differ from the unconstrained errors $E_r$. The distinction strengthens both the general parity hierarchy and the concrete comparison below.

\begin{theorem}[Optimal even-parity staircase]\label{thm:even-staircase}
The numbers $F_r>0$ are attained, nonincreasing, and tend to zero. For fixed $\rho,c,K$, let $\ell=\min\{r:F_r\leq K\}$. The minimum degree of a real even polynomial bounded by one on $[-1,1]$ and approximating $1-x$ within $K\delta$ on \eqref{eq:bands} has the orders in \eqref{eq:staircase}, with this value of $\ell$. In particular,
\begin{equation}\label{eq:even-plateau}
 F_0=G_0,\qquad F_1=F_2=E_1.
\end{equation}
Thus $K=E_1$ admits degree $\Theta(\delta^{-1/2})$, whereas $K<E_1$ forces degree $\Omega(\delta^{-5/6})$.
\end{theorem}
\begin{proof}
A minimizing sequence has bounded values at $2r+1$ fixed distinct points of $[1,\rho]$, hence bounded coefficients. Its coefficientwise limit is even and nonnegative on $\R$, proving attainment. Zero error would identify an even polynomial with $y$ on an interval, an impossibility. Squaring polynomial approximations to $v^{1/4}$ on $[1,\rho^2]$, and then substituting $v=y^2$, proves convergence to zero. Inclusion of the feasible sets proves monotonicity; plateaus are allowed.

For the lower bound with $K<F_r$, repeat the compactness proof of Theorem~\ref{thm:hierarchy} for the rescaled defect $(1-p(\delta y))/\delta$. For each fixed derivative order $j$, interior polynomial Bernstein inequalities give $|p^{(j)}|\leq C_j(d+1)^j$, where $d=\deg p$. Taylor expansion through order $2r+1$ therefore has remainder $O((d+1)^{2r+2}\delta^{2r+1})$ after rescaling, and all its odd coefficients vanish. If $d=o(\delta^{-1+1/(2r+2)})$, a coefficientwise subsequential limit is even, globally nonnegative, and of degree at most $2r$, contradicting $K<F_r$.

For the upper bound let $P$ attain $F_r$, with $r\geq1$. Its positive-error contact set $A=\{y:P(y)-y=F_r\}$ is nonempty: otherwise adding a sufficiently small positive constant would reduce the error while preserving nonnegativity. It is finite because $P(y)-y$ is nonconstant. The polynomial $F_r-P(y)+y$ is nonnegative on $[1,\rho]$; its zeros are exactly $A$, with multiplicities at most $2r$. Thus it is bounded below by $b\dist(y,A)^{2r}$ on that interval for some $b>0$. Lemma~\ref{lem:pinned} and the proof of Proposition~\ref{prop:threshold-upper} now apply with this $P$ and $F_r$ in place of $P_r^*$ and $G_r$: global growth is at most $C(1+|y|^{2r})$, global nonnegativity proves contractivity, and pinning gives exact threshold error. The resulting blend is even because $P$, $W$, and $S_N(x^2)$ are even. Its degree is $O(\delta^{-1+1/(2r)})$. The coarse-tier constructions in Proposition~\ref{prop:coarse} are also even, and its lower bounds apply to these transforms through Lemma~\ref{lem:implementation}. Minimality of $\ell$ gives $K<F_{\ell-1}$ and completes the staircase, including equality cases.

The positive affine polynomial $A_1(v)$ from \eqref{eq:E1} is nonnegative for $v\geq0$, so $F_1=E_1$. To prove $F_2\geq E_1$, write any admissible even quartic as $P(y)=A(y^2)$ with $\deg A\leq2$. Nonnegativity for $v\geq0$ forces its quadratic coefficient to be nonnegative, so $A$ is convex. Let $\ell(v)$ be the chord of $\sqrt v$ on $[1,\rho^2]$, and let the approximation error of $A$ there be $E$. Convexity and the endpoint upper bounds give $A(v)\leq\ell(v)+E$. At $v_*=((\rho+1)/2)^2$, the lower error bound gives $A(v_*)\geq\sqrt{v_*}-E$. Since $\sqrt{v_*}-\ell(v_*)=2E_1$, we have $E\geq E_1$. The admissible affine polynomial already attains equality, proving \eqref{eq:even-plateau}.
\end{proof}

\begin{proposition}[Optimal odd-parity order]\label{prop:odd-optimal}
For fixed $\rho>1$, $0<c\leq1$, and $K>0$, the minimum degree of a real odd polynomial bounded by one on $[-1,1]$ and accurate to $K\delta$ on \eqref{eq:bands} is
\[
 \Theta_{\rho,K}\bigl(\delta^{-1}\log(1/\delta)\bigr).
\]
The lower bound needs only the low interval, and the upper bound holds on all of $[\delta,1]$.
\end{proposition}
\begin{proof}
Let $p$ be such a polynomial of degree $d$. The trigonometric polynomial $q(t)=p(\sin t)$ has degree at most $d$, is bounded by one, and satisfies $q(0)=0$. Set $t_1=\arcsin\delta$, $t_2=\arcsin(\rho\delta)$ and let $A_n$ be its Taylor polynomial at zero through degree $n$. Bernstein's inequality and Taylor's theorem give
\[
 u:=\norm{q-A_n}_{[0,t_2]}
 \leq\frac{(dt_2)^{n+1}}{(n+1)!}.
\]
On $[t_1,t_2]$, $|q(t)-1|\leq(\rho+K)\delta$. The polynomial $1-A_n$ has value one at zero, so exterior Chebyshev gives
\[
 1\leq T_n\left(\frac{t_1+t_2}{t_2-t_1}\right)
            [(\rho+K)\delta+u].
\]
The exterior coordinate tends to $z_0$. Fix any $\overline R>R_0$; for all sufficiently small $\delta$, the Chebyshev factor is at most $\overline R^n$, uniformly in $n$. Take
$n=\lfloor\log(1/\delta)/(2\log\overline R)\rfloor$.
Then $(\rho+K)\delta\overline R^n\to0$, and $u\geq\tfrac12\overline R^{-n}$. Taking $(n+1)$st roots and using the factorial bound yields $dt_2\geq c_\rho n$. Since $t_2=O_\rho(\delta)$, this proves the lower bound.

For the upper bound put $\tau=\min(\delta/2,K\delta/4)$ and $\epsilon=K\delta/4$. For small $\delta$, the bounded odd sign approximant of \cite[Lemma 25, extended version]{gilyen2019} gives $s$ of degree $O(\tau^{-1}\log(1/\epsilon))$, with $|s|\leq1$ on $[-1,1]$ and $|s(x)-1|\leq\epsilon$ for $x\geq\tau$. Define
\[
 p(x)=\frac{s(x)-x}{1+\tau}.
\]
For $0\leq x\leq\tau$, its numerator is at least $-1-\tau$; for $\tau\leq x\leq1$, it is at least $-\epsilon$. Its numerator is at most one on $[0,1]$. Thus $|p|\leq1$ there for small $\delta$, and oddness proves the same bound on $[-1,1]$. For $x\geq\delta$,
\[
 |p(x)-(1-x)|\leq\frac{\epsilon+\tau(1-x)}{1+\tau}
 \leq K\delta.
\]
For fixed $K$, its degree has the stated order.
\end{proof}

For $\rho=2$ the elementary error $\delta/16$ lies above $F_1=1/24$, so Theorem~\ref{thm:even-staircase} gives $\Theta_c(\delta^{-1/2})$ queries even within the even class. At the smaller error $\delta/32$,
\[
 G_1(2)=\frac{3-\sqrt{17/2}}4\approx0.021131
 <\frac1{32}<\frac1{24}=F_1(2)=F_2(2).
\]
To match the even lower bound, set
\[
 z=\frac25y^2-1,\qquad
 P(y)=\sqrt{\frac52}\left(1+\frac z2-\frac{z^2}8+\frac{z^3}{16}\right).
\]
For real $y$, $z\geq-1$. The cubic in parentheses has derivative $1/2-z/4+3z^2/16>0$ and value $5/16$ at $z=-1$, so $P$ is globally positive. For $1\leq y\leq2$, $|z|\leq3/5$. The binomial coefficients of $\sqrt{1+z}$ beyond degree three have absolute values at most $5/128$, and hence
\[
 |P(y)-y|\leq\sqrt{\frac52}\frac5{128}
                    \frac{(3/5)^4}{1-3/5}<\frac1{32}.
\]
Thus $F_3(2)<1/32$ and the even staircase has $\ell=3$. We obtain the following fully matched comparison.
\begin{center}
\begin{tabular}{ll}
\toprule
Converter class & Query order at $\rho=2$, $\eta=\delta/32$\\
\midrule
Arbitrary coherent converter (achieved by GQSP) & $\Theta(\delta^{-1/2})$\\
One even definite-parity QSVT transform & $\Theta(\delta^{-5/6})$\\
One odd definite-parity QSVT transform & $\Theta(\delta^{-1}\log(1/\delta))$\\
\bottomrule
\end{tabular}
\end{center}
The comparison is with a single definite-parity transform, including the usual coherent extraction of its real part. It is not a lower bound on arbitrary compositions of QSVT circuits: such compositions are already allowed in the first row. All these statements count oracle queries and presume exact synthesis of oracle-independent gates; they do not claim a bound on classical phase-computation time or fault-tolerant gate synthesis.
